
\documentclass[12pt,a4paper]{article}

% Page layout and mathematics
\usepackage{geometry}
\usepackage{amsmath}
\usepackage{graphicx}
\usepackage{booktabs}
\usepackage{float}
\usepackage{enumitem}
\usepackage{titlesec}

% Persian typesetting
\usepackage{xepersian}

\geometry{
    top=2.2cm,
    bottom=2.2cm,
    left=2.5cm,
    right=2.5cm
}

\settextfont{B Nazanin}
\setlatintextfont{Times New Roman}

\linespread{1.2}
\setlength{\parskip}{5pt}
\titleformat*{\section}{\large\bfseries}
\titlespacing*{\section}{0pt}{14pt}{6pt}
\widowpenalty=10000
\clubpenalty=10000

\begin{document}

\begin{center}
    {\Large\textbf{پاسخ به مسئله محموله ۱۰۰ هزار دلاری}}
\end{center}

\section*{۱. مسئله}

فرض می‌کنم شرکت متعهد است سه ماه دیگر بابت این محموله ۱۰۰ هزار دلار پرداخت کند و درآمدش ریالی است. پس سؤال
اصلی پیش‌بینی قیمت نیست؛ این است که هزینه ریالی محموله در چه بازه‌ای ممکن است
بیفتد و کدام تصمیم امروز این بازه را محدود می‌کند. هزینه ریالی را این‌طور
تجزیه می‌کنم:

\[
C_{\mathrm{IRR}} = Q\bigl(P_{\mathrm{ICE}}(t_1)+D\bigr)S(t_2) + F + T
\]

که \(Q\) مقدار، \(P_{\mathrm{ICE}}\) قیمت مرجع جهانی (مثلاً فیوچرز عربیکا)،
\(D\) اختلاف قیمت قهوه خریداری‌شده با قیمت مرجع (بسته به کیفیت و کشور مبدأ)، \(S\) نرخ ارز، و \(F\) و \(T\) هزینه حمل و
عوارض است. زمان تثبیت قیمت قهوه (\(t_1\)) و زمان پرداخت ارز (\(t_2\)) لزوماً
یکی نیستند، پس این دو، دو اهرم تصمیم جداگانه‌اند.

\section*{۲. آنچه نمی‌دانم}

پیش از مدل‌سازی این‌ها را از تیم تجاری می‌پرسم:

\begin{itemize}[label=$\circ$]
    \item آیا قیمت با تأمین‌کننده از امروز قابل تثبیت است؟
    \item ارز از چه کانالی و با چه فاصله زمانی تأمین می‌شود؟
    \item موجودی فعلی چند هفته فروش را پوشش می‌دهد؟
    \item اگر دلار گران شود، با چه تأخیری و چند درصد می‌توان قیمت فروش را بالا برد؟
\end{itemize}

سؤال آخر از این جهت مهم است که اگر شرکت بتواند افزایش هزینه را سریع به قیمت
فروش منتقل کند، بخشی از ریسک ارزی عملاً از بین می‌رود. در این حالت ریسک واقعی
کمتر از کل ۱۰۰ هزار دلار است و باید فقط بخشی را اندازه گرفت که به قیمت فروش
منتقل نمی‌شود (\lr{net exposure}). در مقابل، اگر افزایش قیمت فروش با تأخیر
انجام شود، شرکت در این فاصله باید هزینه بالاتر را از نقدینگی خود بپردازد، حتی
اگر این افزایش در نهایت از طریق قیمت فروش جبران شود.

برای کالیبره کردن مدل به این داده‌ها نیاز دارم:

\begin{itemize}[label=$\circ$]
    \item سری تاریخی فیوچرز عربیکا یا روبوستا، بسته به نوع قهوه‌ای که شرکت می‌خرد، و اختلاف قیمت خریدهای قبلی شرکت با آن؛
    \item سری تاریخی نرخ دلار، هم در بازار آزاد و هم در کانالی که شرکت واقعاً با آن ارز تأمین می‌کند؛
    \item فروش، موجودی و قیمت‌های شرکت در دوره‌های قبلی جهش ارزی، برای برآورد سرعت انتقال هزینه به قیمت فروش.
\end{itemize}

\clearpage
\section*{۳. مدل و نتیجه}

برای قهوه از مدل لگ‌نرمال و برای ارز از مدل روند + پرش
(\lr{jump-diffusion}) استفاده کردم، چون ریال روند تضعیف پیوسته دارد به‌علاوه
جهش‌های ناگهانی که توزیع نرمال احتمال آن‌ها را کمتر از واقع برآورد می‌کند. با ۲۰ هزار مسیر
مونت‌کارلو در \lr{Python}، توزیع هزینه ریالی را به‌دست آوردم. پارامترها فرضی‌اند
و فقط برای نمایش روش انتخاب شده‌اند: نوسان سالانه قهوه ۳۵٪، روند تضعیف ریال
۲۵٪ در سال به‌علاوه به‌طور میانگین یک پرش ۱۵ درصدی در سال، و سود سپرده ریالی ۲۸٪. همه اعداد
به‌صورت شاخص گزارش شده‌اند: هزینه ریالی خرید با قیمت قهوه و نرخ ارز امروز برابر ۱۰۰ در نظر گرفته شده است.

\begin{figure}[H]
    \centering
    \includegraphics[width=\linewidth]{2_cost_distribution.png}
\end{figure}

بدون پوشش ریسک، میانگین هزینه ۱۱۱ است، اما در ۵٪ بدترین حالت‌ها هزینه از ۱۵۳
بیشتر می‌شود (\lr{VaR 95\%}) و میانگین همان حالت‌ها ۱۷۰ است
(\lr{Expected Shortfall}). از آنجا که داده‌های تاریخی همه
شوک‌های ممکن را پوشش نمی‌دهند، چند سناریوی مشخص را هم جداگانه بررسی می‌کنم:

\begin{center}
\begin{tabular}{lccc}
\toprule
\textbf{سناریو} & \textbf{نرخ ارز} & \textbf{بدون پوشش} & \textbf{با پوشش*} \\
\midrule
پایه (قهوه بدون تغییر) & روند عادی & ۱۰۶ & ۱۱۴ \\
شوک ارزی & ۳۰٪+ & ۱۳۰ & ۱۱۴ \\
شوک قهوه (۴۰٪+) & روند عادی & ۱۴۹ & ۱۱۴ \\
شوک هم‌زمان (قهوه ۴۰٪+) & ۳۰٪+ & ۱۸۲ & ۱۱۴ \\
کاهش قیمت قهوه (۲۵٪-) & روند عادی & ۸۰ & ۸۸ \\
\bottomrule
\end{tabular}

{\small * خرید کامل ارز از امروز + خرید \lr{\mbox{call option}} قهوه (پرمیوم حدود ۷٪)}
\end{center}

\begin{figure}[H]
    \centering
    \includegraphics[width=\linewidth]{4_strategies.png}
\end{figure}

نتیجه اصلی مقایسه استراتژی‌ها این است که خرید زودهنگام ارز به‌تنهایی دم ریسک را
فقط از ۱۷۰ به ۱۵۳ می‌رساند، چون ریسک قیمت قهوه باز می‌ماند. ترکیب آن با خرید
\lr{\mbox{call option}} قهوه دم ریسک را به ۱۱۴ می‌رساند و هزینه‌اش در حالت عادی حدود
۷٪ است. تثبیت کامل قیمت قهوه از امروز آن را به ۱۰۷ می‌رساند. پس اولویت من تثبیت
قیمت با تأمین‌کننده است، با قرارداد قیمت ثابت (\lr{fixed-price}) یا قرارداد
\lr{PTBF} (\lr{price to be fixed}) که در آن خریدار زمان تثبیت قیمت را انتخاب
می‌کند.

برای ریسک ارز، ابزار مستقیمی مثل فیوچرز یا اختیار معامله دلار-ریال وجود ندارد. جایگزین
قابل بررسی، خرید \lr{\mbox{call option}} روی صندوق‌های طلا در بورس کالاست. قیمت ریالی طلا تقریباً
حاصل‌ضرب قیمت جهانی طلا و نرخ دلار است، پس در جهش‌های ارزی با آن هم‌جهت حرکت
می‌کند. اما این پوشش، ریسک قیمت جهانی طلا و حباب را وارد می‌کند. اینکه برای پوشش ریسک
ارزی این محموله چه مقدار طلا لازم است (نسبت پوشش)، با داده‌های تاریخی تعیین می‌شود. استفاده از این ابزار فقط زمانی قابل توجیه است که کاهش دم ریسک از هزینه
پرمیوم بیشتر باشد.

\section*{۴. حد ریسک و ارزیابی مدل}

حد ریسک را مدیریت تعیین می‌کند و کار من تبدیل آن به قیدهای قابل
اندازه‌گیری است؛ مثلاً «با احتمال ۹۵٪ هزینه از سقفی که نقدینگی و خط اعتباری
پوشش می‌دهد بیشتر نشود» و «احتمال اتمام موجودی پیش از رسیدن محموله زیر ۱۰٪
باشد». از میان استراتژی‌هایی که هر دو قید را رعایت می‌کنند، کم‌هزینه‌ترین را
پیشنهاد می‌کنم.

برای ارزیابی مدل، آن را روی داده‌های گذشته
\lr{backtest} می‌کنم: اگر بازه‌های ۹۵٪ سه‌ماهه در سال‌های گذشته بیش از ۵٪ مواقع
شکسته شده باشند، یعنی مدل احتمال رخدادهای حدی را کمتر از واقع برآورد کرده و باید اصلاح شود. بعد از هر محموله
هم پیش‌بینی را با نتیجه واقعی مقایسه می‌کنم.

\end{document}
